\chapter*{Resumen}

El presente Trabajo Fin de Grado aborda la cuestión central de la \emph{localización de los ceros de los polinomios ortogonales de Sobolev} cuando el producto interior incorpora derivadas de primer orden. Mientras que en la ortogonalidad estándar los ceros permanecen confinados al soporte de la medida, en el contexto Sobolev pueden escapar de dicha región e incluso volverse complejos, fenómeno detectado por Althammer y Meijer que aún no está completamente caracterizado. El objetivo esencial es identificar criterios que garanticen la acotación de los ceros y clarificar la relación entre la geometría de los soportes y la estructura del producto de Sobolev.

Se ha desarrollado un \textbf{marco teórico unificado} que extiende la teoría clásica al caso Sobolev, analizando las propiedades espectrales de las matrices asociadas y conectando la acotación de los ceros con la norma del operador de multiplicación. Esta caracterización pone de relieve el papel diagnóstico de los valores singulares, poco atendidos en la literatura previa.

Se ha implementado en \textsc{Maple} un \textbf{módulo computacional} que genera automáticamente las matrices asociadas, calcula sus propiedades espectrales y visualiza los ceros y soportes en el plano complejo, permitiendo el análisis sistemático de familias de ejemplos.

Con esta infraestructura se han analizado tres \textbf{familias representativas}:

\emph{Caso real}: Para medidas reales con derivadas de primer orden, los experimentos confirman que los ceros permanecen acotados cuando el operador es acotado, convergiendo al radio de la envoltura convexa.

\emph{Caso circunferencia}: Se han reproducido y extendido ejemplos de circunferencias concéntricas, tangentes y disjuntas, confirmando que la acotación depende de la envoltura convexa global y no de la conectividad del soporte.

\emph{Caso discreto}: En la mezcla de una medida continua con un punto aislado se observa la pérdida repentina de acotación cuando algunos ceros escapan hacia el punto aislado, evidenciando la necesidad de condiciones adicionales para asegurar estabilidad.

De la combinación de teoría y experimentación surgen dos \textbf{conjeturas}: (i) en todo producto Sobolev con soportes compactos, el límite del radio espectral coincide con el radio máximo de la envoltura convexa combinada; (ii) cuando el operador no es acotado, el segundo valor singular converge al máximo módulo de los puntos de soporte.

Las \textbf{aportaciones principales} incluyen: (1) unificar criterios de acotación dispersos en la bibliografía; (2) introducir herramientas gráficas para diagnosticar la acotación; (3) aportar evidencia sobre la influencia de puntos aislados y soportes no conexos; y (4) formular conjeturas que refinarían las cotas espectrales actuales.

El trabajo ilustra la sinergia entre análisis teórico y experimentación numérica, mostrando cómo la visualización puede guiar la intuición matemática. Sus resultados consolidan teoremas clásicos, los extienden a geometrías más generales y abren líneas futuras hacia productos de orden superior y medidas matriciales.

\textbf{Palabras clave:} Polinomios ortogonales de Sobolev, localización de ceros, análisis espectral, visualización numérica.


\begin{otherlanguage}{english}
  \section*{Abstract}
  This Final Degree Project tackles the central question of the \emph{location of the zeros of Sobolev orthogonal polynomials} when the inner product includes first-order derivatives.  Whereas in standard orthogonality the zeros remain confined to the support of the measure, in the Sobolev setting they may escape that region and even become complex—a phenomenon detected by Althammer and Meijer that is still not fully characterised.  The essential goal of the work is to identify criteria that guarantee the boundedness of the zeros and to clarify the relationship between the geometry of the supports and the structure of the Sobolev product.

  To this end a \textbf{unified theoretical framework} has been developed, extending the classical theory to the Sobolev case, analysing the spectral properties of the associated matrices, and connecting the boundedness of the zeros with the norm of the multiplication operator. This characterisation highlights the diagnostic role of the singular values, scarcely addressed in previous literature.

  On this basis, a \textbf{computational module} has been implemented in \textsc{Maple} that automatically generates the associated matrices, computes their spectral properties, and visualises the zeros and supports in the complex plane, allowing for the systematic analysis of families of examples.

  With this infrastructure, three \textbf{representative families} have been analysed:

  \emph{Real case}: For real measures with first-order derivatives, experiments confirm that the zeros remain bounded when the operator is bounded, converging to the radius of the convex hull.

  \emph{Circular case}: The examples of concentric, tangent, and disjoint circles have been reproduced and extended, confirming that the boundedness depends on the global convex hull and not on the connectivity of the support.

  \emph{Discrete case}: In the mixture of a continuous measure with an isolated point, a sudden loss of boundedness is observed when some zeros escape towards the isolated point, evidencing the need for additional conditions to ensure stability.

  From the combination of theory and experimentation, two \textbf{conjectures} emerge: (i) for every Sobolev product with compact supports, the limit of the spectral radius coincides with the maximum radius of the combined convex hull; (ii) when the operator is unbounded, the second singular value converges to the maximum modulus of the support points.

  The \textbf{main contributions} include: (1) unifying scattered criteria for boundedness in the literature; (2) introducing graphical tools to diagnose boundedness; (3) providing evidence on the influence of isolated points and non-connected supports; and (4) formulating conjectures that would refine current spectral bounds.

  Overall, the work illustrates the \emph{synergy between theoretical analysis and numerical experimentation}, showing how visualisation can guide mathematical intuition. Its results consolidate classical theorems, extend them to more general geometries, and open future lines towards higher-order products and matrix measures.

  \textbf{Keywords:} Sobolev orthogonal polynomials, location of zeros, spectral analysis, numerical visualisation.

\end{otherlanguage}


%%%%%%%%%%%%%%%%%%%%%%%%%%%%%%%%%%%%%%%%%%%%%%%%%%%%%%%%%%%
%% Final del resumen.
%%%%%%%%%%%%%%%%%%%%%%%%%%%%%%%%%%%%%%%%%%%%%%%%%%%%%%%%%%%